\chapter{Introduction}

The increasing digitisation of financial services has created a growing need for
privacy-preserving collaborative computation. Financial institutions must detect fraud
and screen customers against regulatory watchlists maintained by bodies such as the
Office of Foreign Assets Control (OFAC), the Financial Action Task Force (FATF), and
the UN Security Council --- yet privacy regulations such as the European General Data
Protection Regulation (GDPR) prohibit sharing raw customer data with third parties.
Cryptographic Private Set Intersection (PSI) addresses this tension by allowing two
parties to jointly compute the intersection of their datasets without revealing any
element outside the intersection.

A second, longer-term concern is quantum computing. Most deployed PSI protocols rely
on Diffie-Hellman or elliptic-curve assumptions that can be broken in polynomial time
by a quantum computer running Shor's algorithm. Given that sensitive financial records
carry regulatory retention obligations of 7--30 years, data encrypted today under
quantum-vulnerable cryptography is already at risk of future decryption by adversaries
who store ciphertexts now and decrypt later --- the so-called \textit{harvest-now,
decrypt-later} (HNDL) attack. This motivates building PSI on post-quantum foundations.

% ─────────────────────────────────────────────────────────────────────────────
\section{Background}

\begin{definition}[Private Set Intersection]\label{def:psi}
Let $\mathcal{P}_1$ hold a private set $S \subseteq \mathcal{U}$ of size $n$ and
$\mathcal{P}_2$ hold a private set $C \subseteq \mathcal{U}$ of size $m$, over a
common universe $\mathcal{U}$. A Private Set Intersection (PSI) protocol is a
two-party protocol that allows both parties to jointly compute $S \cap C$ such that:
\begin{itemize}
    \item $\mathcal{P}_1$ (the server) learns nothing beyond $|C|$.
    \item $\mathcal{P}_2$ (the client) learns nothing beyond $S \cap C$.
\end{itemize}
\end{definition}

\begin{definition}[Laconic PSI]\label{def:laconic}
A PSI protocol is \textit{laconic} if the server's first message has length
$O(\lambda \cdot \log n)$, independent of the server set size $n$, where $\lambda$
is the security parameter~\cite{dottling2023efficient}.
\end{definition}

\begin{remark}
The laconic property is particularly valuable when a single server communicates with
many clients over time. The server publishes one short digest (the Merkle root of
its dataset), and each client sends $O(m)$ ciphertexts in a second message.
Communication thus scales with the client set size, not the server set size. This
asymmetry is precisely what makes laconic PSI attractive for sanctions screening,
where the server holds a large, slowly-changing list queried by many institutions.
\end{remark}

\begin{definition}[Post-Quantum Security]\label{def:pq}
A cryptographic scheme is post-quantum secure if the best known algorithm for
breaking its underlying hardness assumption runs in time $2^{\Omega(\lambda)}$ even
on a quantum computer, where $\lambda$ is the security parameter. Quantum computers
running Shor's algorithm can break Diffie-Hellman and discrete-logarithm-based
assumptions in polynomial time, but no efficient quantum algorithm is known for
lattice problems such as Learning With Errors (LWE)~\cite{regev2005lattices}.
\end{definition}

% ─────────────────────────────────────────────────────────────────────────────
\section{Motivation}

Financial institutions are legally required to screen customers against sanctions
lists to prevent money laundering and terrorism financing. This creates a precise
privacy conflict:
\begin{itemize}
    \item The bank must determine whether customer $x$ appears on the sanctions list.
    \item The bank must not reveal the identity of $x$ to the list authority.
    \item The list authority must not reveal the full sanctions list to the bank.
\end{itemize}

Classical PSI protocols such as KKRT16~\cite{kolesnikov2016efficient} address this
efficiently, but their base Oblivious Transfer is typically instantiated with
elliptic-curve cryptography, which is vulnerable to quantum attacks. While it is
theoretically possible to replace the base OT with a lattice-based construction
(making the full KKRT pipeline potentially post-quantum), no such instantiation has
been standardised or benchmarked at scale. Furthermore, KKRT and similar OT-based
protocols achieve $O(n + m)$ communication complexity, which grows linearly with the
server set size $m$ --- a significant drawback when $m$ is large (e.g., a global
sanctions database of millions of entries).

Existing post-quantum PSI alternatives either incur $O(n + m)$ communication cost
or exist only as theoretical constructions without empirical validation at
real-world dataset sizes. This project directly addresses this gap by implementing
and empirically evaluating the first post-quantum laconic PSI system.

% ─────────────────────────────────────────────────────────────────────────────
\section{Problem Statement}

This project addresses the following central problem:
\begin{quote}
    \textit{How can we implement a post-quantum-secure, laconic Private Set Intersection
    protocol that scales to real-world financial dataset sizes on commodity hardware,
    while maintaining cryptographic correctness and GDPR compliance?}
\end{quote}

The challenge is two-dimensional. First, the memory requirement of the D\"{o}ttling
et al.\ (DKLLMR23) protocol~\cite{dottling2023efficient} presents a fundamental
scalability barrier. In the na\"{i}ve fully-parallel implementation, the server must
hold a witness vector for every element simultaneously, causing auxiliary space to
grow as:
\[
    S_{\text{aux}}(n) = O(n \cdot D \cdot \log q)
\]
where $D$ is the ring dimension and $q$ is the modulus. At practical dataset sizes,
this evaluates to hundreds of gigabytes --- far exceeding commodity server capacity.
A na\"{i}ve single-threaded sequential implementation avoids this by reusing a single
witness buffer, but at the cost of being proportionally slower with dataset size.
Neither extreme is acceptable for deployment.

Second, even with memory resolved, the computational cost of Ring-LWE operations means
that single-node execution remains too slow for time-sensitive compliance workflows.
A distributed execution strategy is therefore required to bring latency within
operationally acceptable bounds.

This project addresses both dimensions: the memory problem through batched witness
generation, and the latency problem through distributed parallel processing across a
cloud node cluster.

% ─────────────────────────────────────────────────────────────────────────────
\section{Objectives}

This project has the following concrete objectives, corresponding to work done in
Phase~I and Phase~II:

\paragraph{Phase I (completed):}
\begin{itemize}
    \item Study and formalise the DKLLMR23 laconic PSI protocol.
    \item Implement the full Ring-LWE-based Laconic Encryption primitive (KeyGen,
          Enc, Dec, WitGen) in Go.
    \item Encrypt financial transaction records and store them in a SQLite-backed
          Merkle tree.
    \item Validate correctness on small datasets ($n \le 250$).
\end{itemize}

\paragraph{Phase II (this report):}
\begin{itemize}
    \item Resolve the memory bottleneck via batched witness generation, reducing
          auxiliary space from $O(n \cdot D \cdot \log q)$ to $O(B \cdot D \cdot \log q)$
          where $B \ll n$ is the batch size.
    \item Implement an in-memory Merkle tree, eliminating database I/O during
          witness generation ($21\times$ speedup).
    \item Validate scalability to $n = 10{,}000$ at 64-bit quantum security ($D = 256$)
          on an AMD EPYC HPC server.
    \item Validate cryptographic correctness at 128-bit post-quantum security ($D = 2048$).
    \item Design and deploy a distributed execution architecture across a Google Compute
          Engine cluster of 20--25 nodes, reducing end-to-end execution time from 153
          minutes (single node) to approximately 10 minutes --- achieving a $15\times$
          speedup through horizontal partitioning of the server dataset and parallel
          intersection detection across nodes.
    \item Develop FLARE, a GDPR-compliant proof-of-concept sanctions screening system.
\end{itemize}

% ─────────────────────────────────────────────────────────────────────────────
\section{Contributions Summary}

To clearly delineate the scope of this project relative to prior work, the
contributions are summarised as follows:

\begin{enumerate}
    \item \textbf{First open implementation of DKLLMR23.} No prior open-source
          implementation or empirical benchmark of this protocol existed.
    \item \textbf{Batched witness generation.} A memory management scheme that
          reduces peak RAM from $\approx 312$\,GB to 6.5\,GB (97.9\% reduction)
          while preserving parallel speedup.
    \item \textbf{In-memory Merkle tree.} Eliminating SQLite I/O during witness
          generation yields a $21\times$ speedup in that phase.
    \item \textbf{Distributed execution on Google Compute Engine.} Horizontal
          partitioning of the server dataset across a 20--25 node GCE cluster
          reduces end-to-end execution time from 153 minutes to approximately 10
          minutes --- a $15\times$ speedup --- making the protocol viable for
          time-sensitive overnight compliance workflows.
    \item \textbf{Empirical benchmark.} The first measured performance data for
          Ring-LWE-based laconic PSI, complementing the pairing-based empirical
          study of ALOS22.
    \item \textbf{FLARE prototype.} A complete, running GDPR-compliant sanctions
          screening web application demonstrating real-world viability.
\end{enumerate}

% ─────────────────────────────────────────────────────────────────────────────
\section{Organisation of the Report}

Chapter 2 surveys related work in PSI and laconic cryptography, incorporating
corrections to prior post-quantum status claims based on reviewer feedback.
Chapter 3 presents the theoretical background including formal complexity analysis.
Chapter 4 describes the system architecture and protocol design.
Chapter 5 details the FLARE prototype implementation.
Chapter 6 presents and analyses experimental results.
Chapter 7 concludes with an honest discussion of limitations and future directions.

\clearpage
